%-------------------------------------------------------------------------------
%	ХУРААНГУЙ
%-------------------------------------------------------------------------------

\addchaptertocentry{Хураангуй}

\begin{center}
{\scshape\Large \univname\par}
{\scshape\large \facname\par}\vspace{0.5cm}
{\huge\textbf{{Хураангуй}} \par}
\bigskip
{\Large{\ttitle} \par}
\bigskip

{\normalsize \shortname \par}
\addressname
\end{center}

\textit{\textbf{Түлхүүр үгс: \keywordnames}}
\bigskip

Санхүүгийн байгууллагын ажилтнууд өгөгдлийн сангаас тоо мэдээ авахдаа SQL бичдэг шинжээчээс хамаардаг. Байгалийн
хэлний асуултыг SQL болгох (Text2SQL) системүүд англи хэлний Spider, BIRD жишиг дээр өндөр үр дүн үзүүлж байгаа ч
монгол хэлэнд зориулсан систем, үнэлгээний сан, туршилтын нотолгоо бараг байхгүй. Монгол хэлний залгамал үгзүй,
кирилл ба латин галигийн холимог бичлэг, өгөгдлийн санд монголоор хадгалагдсан утга, байгууллагын дотоод нэр томьёо
(PAR30, идэвхтэй багц) нь англи хэлэнд боловсруулсан аргуудыг шууд хэрэглэхэд саад болдог. Мөн өгөгдөл гадагш гарах
ёсгүй тул нэг GPU-д багтах нээлттэй эхийн загвар ашиглах шаардлагатай.

Энэхүү ажлаар банк бус санхүүгийн байгууллагын (ББСБ) 28 хүснэгт бүхий PostgreSQL өгөгдлийн сан (648 205 мөр
синтетик өгөгдөл) дээр монгол асуултад хариулах Text2SQL системийг нээлттэй эхийн Arctic-Text2SQL-R1-7B загвар дээр
бүтээж, бүрэлдэхүүн бүрийн нөлөөг хяналттай туршилтаар хэмжив. Үүний тулд dev 300, test 200 асуулт бүхий монгол
үнэлгээний сан (хариулах боломжгүй ба хоёрдмол асуулттай) бэлтгэж, загварын гэр бүлээр нь давхцалгүй хуваав. Систем нь
асуултыг шалгах (хувийн мэдээлэл, бичих хүсэлтээс татгалзах), Qwen3.5-9B-ээр англи руу орчуулах (утгыг маскаар
хамгаалах), мэдлэг хайх (BM25 ба BGE-M3: хүснэгт, утга, нэр томьёо, жишээ), өөрийн DB-ийн RAG prompt-оор LoRA-аар
нэмж сургасан загвараар SQL үүсгэх, алдааг засах, 8 хувилбараас гүйцэтгэлийн үр дүнгээр санал хураах, итгэл багатай
үед тодруулга асуух алхмуудаас бүрдэнэ. Мөн туршилтын тохиргоог node-оор угсарч, алхам бүрийн prompt ба мэдлэгийг
харуулдаг «Text2SQL Studio» вэб хэрэгслийг хөгжүүлэв.

Dev дээр нийт зөв хариулт суурь загварын 2.3\%-иас (RAG-гүй) 53.0\% хүртэл өсөв. Шинэ асуултаас бүрдсэн test дээр
нэмэлт сургалттай суурь (D2) 21.0\% байсан бол орчуулга ба засвартай H4 34.0\% (McNemar $p = 8.6 \times 10^{-4}$),
8 хувилбарын санал хураалттай E2 38.5\% (хариулах боломжтой асуултын EX 41.1\%), тодруулгатай E3c 46.5\%, жишээний
сан ба орчуулгын засвартай I3tr 45.0\% болов. Бүрэлдэхүүнээр нь задлахад мэдлэг хайх (+12.0 пункт) ба орчуулга (+11.3)
хамгийн их, өөрийн DB-ийн жишээгүй нэмэлт сургалт (+1.4) хамгийн бага нэмэр өгсөн. Алдааны шинжилгээгээр E2-ийн
алдааны 38\% нь орчуулгаас үүдсэн тул санхүүгийн хэллэгийн 2 828 хос бүхий орчуулгын өгөгдөл ба схем, тодорхойлолт
өөрчлөгдөхөд тогтвортой байх SQL бичигчийн сургалтын өгөгдлийг бэлтгэв. Ажлын гол хувь нэмэр нь монгол хэлний анхны
Text2SQL үнэлгээний сан, бүрэлдэхүүн тус бүрийн нөлөөг статистик тестээр баталсан туршилт, нэг GPU дээр ажиллах
давтагдахуйц систем болон судалгааны хэрэгсэл юм.
